%==============================================================
% PRAKATA
%==============================================================

\chapter*{PRAKATA}
\addcontentsline{toc}{chapter}{PRAKATA}

Puji dan syukur senantiasa dipanjatkan ke hadirat Tuhan Yang Maha Esa atas
segala karunia-Nya sehingga laporan project akhir yang berjudul
\textit{Klasifikasi Stage Retinopathy of Prematurity Berdasarkan Citra Fundus Bayi
Menggunakan Representasi Spasial Softmap dan TinyResNet}
dapat diselesaikan dengan baik. Laporan ini disusun sebagai pemenuhan syarat
penilaian akhir pada mata kuliah Pengenalan Citra Digital.

Penyelesaian laporan ini tidak lepas dari bimbingan dan bantuan berbagai pihak.
Oleh karena itu, penulis mengucapkan terima kasih kepada dosen pengampu mata
kuliah Pengenalan Citra Digital yang telah memberikan arahan teknis maupun
konseptual selama pengerjaan project. Penulis juga mengapresiasi kerja keras
seluruh anggota kelompok yang telah berkontribusi aktif sejak tahap eksplorasi
dataset, implementasi pipeline, hingga penyusunan laporan ini.

Penulis menyadari bahwa masih terdapat ruang pengembangan, terutama pada
representasi area tepi yang membedakan Stage 1 dan Stage 2 serta pengujian model
pada data klinis yang lebih beragam. Penulis berharap laporan ini dapat menjadi
landasan yang baik untuk penelitian lanjutan mengenai klasifikasi ROP berbasis
citra fundus bayi.

\vspace{1.5cm}

\begin{flushright}
Bogor, Juni 2026\\[1.5cm]
Kelompok 15
\end{flushright}
